% Source MD: E:\wanru\PRML_Course\NoteMD\Week1\第2节课 深度学习基本概念.md
\documentclass[main.tex]{subfiles}

\graphicspath{{\subfix{figures/}}}

\begin{document}

\chapter{深度学习基本概念}
\label{chap:w01l02-dl-intro}

\begin{intuitionbox}
深度学习没有突然换一套新游戏规则。它做的事比较像是在原来的线性模型上加更多弹性，让函数弯得起来，然后继续沿用同一套 loss 和优化流程训练参数。
\end{intuitionbox}

\section{线性模型的局限与模型偏差}
\label{sec:w01l02-model-bias}

先从上一节的线性模型继续往下看。$y=b+wx_1$ 这种写法当然简单，但它有一个很硬的限制：不管 $w$ 和 $b$ 怎么调，最后都只能画出一条直线。
如果真实关系本来就弯弯折折，这种模型再怎么训，也很难贴上去。

\begin{definition}{模型偏差（Model Bias）}{w01l02-model-bias}
如果模型本身能表示的函数集合太小，就算参数已经调到这类模型里最好的状态，仍然逼不近真实关系，这种由表达能力不足带来的误差，就叫模型偏差（model bias）。
\end{definition}

所以这节的重点不是“把线性模型调得更努力”，而是先承认它不够，再换一个更有弹性的函数家族。

\section{分段线性曲线与 Sigmoid}
\label{sec:w01l02-sigmoid}

课程一开始给的直觉其实很朴素：一条复杂曲线，不一定要一次写完，也可以想成很多小段慢慢拼出来。
只要转折点够多，分段线性曲线（piecewise linear curve）就已经能逼近很多连续函数。

问题在于，真的把每一段折线都手写出来会很麻烦。于是课程换了一个更好调的零件，也就是 sigmoid：
\begin{equation}
  y = c \cdot \frac{1}{1+e^{-(b+wx_1)}} = c \cdot \sigmoid\paren{b+wx_1}.
  \label{eq:w01l02-sigmoid}
\end{equation}

\cref{eq:w01l02-sigmoid} 真正控制形状的是中间那项 $z=b+wx_1$。
当 $z$ 很小的时候，输出会压在 0 附近；当 $z$ 很大的时候，输出又会慢慢靠近 $c$。
也就是说，它不像直线那样一路往前冲，而是会出现一个从低平台过渡到高平台的区段。用它来模仿 hard sigmoid，就会比硬折线顺手很多。

$w$ 决定这段过渡是陡还是缓，$b$ 决定转折大概出现在横轴哪里，$c$ 则把整段输出的高度拉到目标尺度。
先把这三个角色分清，后面看多个 sigmoid 叠加时就不会乱。

\begin{examplebox}
如果把 sigmoid 的转折区想成“开关开始打开的那一段”，那它大致就落在 $b+wx_1=0$ 附近，也就是 $x_1 \approx -b/w$。
这个位置会随着 $b$ 左右移动，过渡有多突然则更受 $w$ 影响，至于最后平台能抬多高，就看 $c$。
\end{examplebox}

单个 sigmoid 当然还是有限，但这类元件的好处就在于可以堆。
一条不够，就放很多条；一个人弯不过来，就让一群元件一起分工。深度学习里很常见的做法，本质上就是把许多可调的非线性单元拼起来，让模型自己学会怎么组合。

\section{更有弹性的模型}
\label{sec:w01l02-flexible-model}

沿着这个思路往前走，把多个 sigmoid 叠起来，模型就会开始有明显的弹性：
\begin{equation}
  y = b + \sum_i c_i \cdot \sigmoid\paren{b_i + w_i x_1}.
  \label{eq:w01l02-single-feature}
\end{equation}

\cref{eq:w01l02-single-feature} 里那个索引 $i$，就是在数“第几个 sigmoid 单元”，也可以把它理解成第 $i$ 个隐藏神经元（hidden neuron）。
每个单元都有自己的一套 $w_i,b_i,c_i$：前两者决定这条小曲线在哪里转、转得有多快，后者则决定它对最终输出贡献多大。
许多小曲线叠在一起，模型就不再只能画一条直线。

如果输入不只一个特征，而是 $x_1,x_2,\dots$ 一整组，写法就会自然扩成
\begin{equation}
  y = b + \sum_i c_i \cdot \sigmoid\paren{b_i + \sum_j w_{ij}x_j}.
  \label{eq:w01l02-multi-feature}
\end{equation}

$j$ 现在用来数输入特征，$w_{ij}$ 则表示“第 $j$ 个特征对第 $i$ 个神经元有多大影响”。
所以对某个固定的神经元来说，它看到的不是单独某一维输入，而是一整组特征先加权混合后的结果 $\sum_j w_{ij}x_j$。
也正因为每个神经元都能从不同角度去看同一批输入，多个 sigmoid 叠起来以后，模型的可塑性才会明显上来。

\begin{notation}{参数向量 $\theta$}{w01l02-theta}
把矩阵 $W$、偏置向量 $b$、输出层权重 $c$ 以及最终常数项全部拉直并接在一起，就得到参数向量 $\theta$。它代表模型里所有需要学习的未知数。
\end{notation}

\subsection{矩阵表示}
\label{sec:w01l02-matrix-form}

上面的标量写法已经够说明问题，但一旦神经元和特征数量变多，继续一个个写下标会很快失控。
所以后面通常会把同样的计算压缩成矩阵形式：
\begin{align}
  r &= Wx + b, \label{eq:w01l02-r} \\
  a &= \sigmoid\paren{r}, \label{eq:w01l02-a} \\
  y &= b_{\mathrm{final}} + c\trans a. \label{eq:w01l02-y}
\end{align}

\cref{eq:w01l02-r,eq:w01l02-a,eq:w01l02-y} 没有换掉原本的逻辑，只是把“每个神经元各算一次”压缩成一次向量化运算。
设输入向量 $x \in \R^d$，隐藏层共有 $m$ 个神经元（neuron），那么：
\begin{itemize}
  \item $W \in \R^{m \times d}$：第 $i$ 行收集了第 $i$ 个神经元对所有特征的权重；
  \item $b \in \R^m$：第 $i$ 个分量对应第 $i$ 个神经元的偏置；
  \item $r \in \R^m$：其中第 $i$ 个分量就是原来那个 $r_i=\sum_j w_{ij}x_j+b_i$；
  \item $a \in \R^m$：第 $i$ 个分量满足 $a_i=\sigmoid(r_i)$；
  \item $c \in \R^m$：拿来把所有隐藏层输出重新加权汇总。
\end{itemize}

矩阵写法的好处不只是省篇幅，它也把神经网络最常见的计算顺序露得很清楚：先做仿射变换，再套 activation function，最后把这一层的输出送到下一层。

\begin{examplebox}
如果输入只有两个特征，隐藏层有三个神经元，那么 $x$ 就是 $2$ 维向量，$W$ 会是一个 $3 \times 2$ 的矩阵。
这表示三个神经元都在看同一组输入，但它们看的角度不一样，所以学出来的非线性模式也会不同。
\end{examplebox}

\section{训练：Loss、优化、Batch 与 Epoch}
\label{sec:w01l02-training}

模型开始变复杂以后，训练流程本身其实没换。
前面多出来的是函数表达能力，不是训练的基本逻辑。照样还是先写 loss，再去解
\begin{equation}
  \theta^\star = \argmin_\theta L\paren{\theta}.
  \label{eq:w01l02-opt}
\end{equation}

更新规则也还是熟悉的梯度下降：
\begin{equation}
  \theta^{t+1} = \theta^t - \eta \grad L\paren{\theta^t}.
  \label{eq:w01l02-gd}
\end{equation}

\cref{eq:w01l02-opt,eq:w01l02-gd} 看起来抽象一些，只是因为参数太多，不能再像第一课那样逐个写 $w,b$。
这里的 $\theta$ 把 $W,b,c,b_{\mathrm{final}}$ 全部包起来，所以梯度 $\grad L(\theta^t)$ 也会是一条同样长度的向量，里面每个分量都对应着某个参数该往哪里调。

\begin{definitionbox}[title={Batch 与 Epoch}]
实际训练时，通常不会每次都用完整训练集来算一次梯度，而是把数据切成许多 batch。
看完一个 batch 并更新一次参数，称为一次 update；把所有 batch 全看完一遍，称为一个 epoch。
\end{definitionbox}

换成 batch 的角度去看，就是每次只拿一部分资料来估计完整 loss 的下降方向。
单次梯度没那么精确，但计算便宜得多，参数也能更频繁更新，这就是为什么实际训练几乎都会这么做。

\section{Activation Function 与 Deep Learning}
\label{sec:w01l02-activation}

如果说 sigmoid 像一种比较平滑的非线性元件，那 ReLU 就是更直接的另一种版本：
\begin{equation}
  y = c \cdot \max\paren{0, b+wx_1}.
  \label{eq:w01l02-relu}
\end{equation}

两者都能把“纯直线”这件事打破，不过手法不一样。
sigmoid 是平滑地从低平台过渡到高平台；ReLU 更干脆，负的部分直接砍成 0，正的部分保留成直线。
所以前者像平滑开关，后者更像硬一点的折线。

sigmoid 和 ReLU 都属于激活函数（activation function）。
一层算完再接一层，模型就会往深的方向长。课程里的命名关系可以顺手整理成：
\begin{itemize}
  \item 一个 sigmoid 或 ReLU 单元可以视为一个神经元（neuron）；
  \item 一排神经元组成一个隐藏层（hidden layer）；
  \item 隐藏层堆得比较深时，这整套方法就会被叫做深度学习（deep learning）。
\end{itemize}

\begin{examplebox}
课程实验里可以很直接地看到：神经元数量和层数一上去，训练误差常常就会往下掉。
不过模型一旦太会拟合，后面也就更容易接到过拟合的问题。
\end{examplebox}

\section{过拟合现象}
\label{sec:w01l02-overfitting}

模型变强当然是好事，但也不是没有代价。
课程还是用 YouTube 观看人数预测举例：更深的网络在训练集上可能压出更低的 loss，可一到 2021 年这种没见过的资料上，表现反而变差。

\begin{definition}{过拟合（Overfitting）}{w01l02-overfitting}
如果模型在训练资料上越学越好，但换到没见过的测试资料时反而变差，这种现象就叫过拟合（overfitting）。
\end{definition}

所以选模型时，不能只盯着训练误差看。真正该看的，是它遇到未知资料时还有没有泛化能力。

\begin{summarybox}
这节的主线可以顺成一句话：线性模型不够弯，所以我们引入非线性元件；元件一多，模型就会变成神经网络；层数继续往上叠，深度学习也就自然出现了。
\end{summarybox}

\end{document}
